% !TEX program = xelatex
% 理论力学 · 专题讲义 · 从张量到并矢
\documentclass[11pt]{ctexart}
%=====================================================================
% theme-violet.tex —— 讲义精讲 · 紫花主题
% 用法：\documentclass[11pt]{ctexart}
%       \input{assets/theme-violet}
%       \begin{document} ... \end{document}
% 只用 ctexart 自带能力 + xcolor/enumitem/fancyhdr/pifont/amsmath，
% 不依赖 tikz/tcolorbox，保证 MiKTeX 可直接编译。
%=====================================================================

\usepackage[a4paper,left=2.0cm,right=2.0cm,top=2.1cm,bottom=2.3cm,%
            headsep=12pt,footskip=22pt]{geometry}
\usepackage{amsmath,amssymb}
\usepackage{needspace}
\usepackage{xcolor}
\usepackage{array}
\usepackage{enumitem}
\usepackage{pifont}
\usepackage{fancyhdr}
\usepackage[hidelinks]{hyperref}

%---------------------------------------------------------------- 配色
\definecolor{violet}{RGB}{107,78,143}      % 主色：紫
\definecolor{violetdeep}{RGB}{63,44,92}    % 深紫：标题/强调
\definecolor{petal}{RGB}{183,159,208}      % 花瓣：装饰线、花饰
\definecolor{bandviolet}{RGB}{247,243,251} % 浅紫底：档案/提示框
\definecolor{bandcream}{RGB}{252,249,242}  % 米白底：易错/注意框
\definecolor{rulepurple}{RGB}{217,207,230} % 分隔线
\definecolor{inkgray}{RGB}{92,96,110}      % 次要文字

%---------------------------------------------------------------- 字体
% 有 Noto CJK 用 Noto，没有就退回 ctex 默认（Windows 下为宋体/黑体）
\IfFontExistsTF{Noto Serif CJK SC}{\setCJKmainfont{Noto Serif CJK SC}}{}
\IfFontExistsTF{Noto Sans CJK SC}{\setCJKsansfont{Noto Sans CJK SC}}{}

%---------------------------------------------------------------- 版式
\setlist{nosep,leftmargin=1.9em}
\linespread{1.28}
\setlength{\parskip}{3.5pt}
\setlength{\parindent}{0pt}
\setlength{\tabcolsep}{5pt}
\renewcommand{\arraystretch}{1.28}

%---------------------------------------------------------------- 花饰
% \ding{95} = 五瓣小花，\ding{96} = 星形花；都来自 pifont，必定存在
\newcommand{\orn}{\textcolor{petal}{\ding{95}}}          % 行内小标
\newcommand{\ornstar}{\textcolor{petal}{\ding{96}}}      % 星形花
\newcommand{\flowerrule}{%                                % 居中花饰分隔线
  \par\vspace{7pt}%
  \noindent\makebox[\linewidth]{\color{petal}%
    \rule{0.20\linewidth}{0.6pt}\hspace{8pt}{\small\ding{95}}\hspace{4pt}%
    \raisebox{1.1pt}{\tiny\ding{95}}\hspace{4pt}{\small\ding{95}}\hspace{8pt}%
    \rule{0.20\linewidth}{0.6pt}}%
  \par\vspace{7pt}}

%---------------------------------------------------------------- 标题与节头
\newcommand{\doctitle}[3]{%   #1 主标题  #2 副标题  #3 元信息（可用 \\ 分行）
  \begin{center}
    {\color{petal}\small \orn\enspace #2\enspace \ornstar}\\[7pt]
    {\LARGE\bfseries\color{violetdeep} #1}\\[10pt]
    {\color{petal}\rule{0.55\linewidth}{1pt}}\\[4pt]
    {\small\color{inkgray} #3}
  \end{center}\vspace{2pt}}

\newcommand{\hdr}[2]{%   #1 章节号（如 §0）  #2 标题
  \par\vspace{12pt}%
  \noindent\colorbox{violetdeep}{\color{white}\bfseries\small\;#1\;}\hspace{7pt}%
  {\color{violetdeep}\large\bfseries #2}\par\nobreak
  \vspace{3pt}{\color{petal}\hrule height 1.2pt}\vspace{6pt}}

\newcommand{\subhdr}[1]{%
  \par\vspace{6pt}%
  \noindent{\color{violet}\bfseries #1}\par\vspace{1pt}}

%---------------------------------------------------------------- 提示框
\newcommand{\callout}[1]{% 浅紫底：讲义档案、读法说明
  \par\medskip\noindent
  \fcolorbox{rulepurple}{bandviolet}{\parbox{\dimexpr\linewidth-2\fboxsep-2\fboxrule}{\vspace{2pt}#1\vspace{2pt}}}%
  \par\medskip}

\newcommand{\warnbox}[1]{% 米白底：易错、边界条件
  \par\medskip\noindent
  \fcolorbox{rulepurple}{bandcream}{\parbox{\dimexpr\linewidth-2\fboxsep-2\fboxrule}{\vspace{2pt}#1\vspace{2pt}}}%
  \par\medskip}

\newcommand{\gapitem}[2]{% #1 缺口类型（紫） #2 内容
  \par\vspace{2pt}\noindent
  \hangindent=1.4em\hangafter=1
  {\color{petal}\small\ding{96}}\enspace{\bfseries\color{violetdeep}#1}\enspace #2\par}

% 题干块（\qstem）：每道例题／作业的「详解」开头必须先复述题干与出处。
% 起因：2026-10-25 第五章讲义把七道作业的题面集中放在一页、详解另起一页，
% 读者翻到「第 4 题」时看不到条件，等于一道没有题干的题目。
% 规则：详解块内自带题干，读者不翻页就能读懂题在问什么。
\newcommand{\qstem}[1]{%
  \par\medskip\noindent
  \fcolorbox{violet}{bandviolet}{%
    \parbox{\dimexpr\linewidth-2\fboxsep-2\fboxrule}{%
      \vspace{2pt}%
      {\small\bfseries\color{violet}题干}\hspace{0.8em}#1%
      \vspace{2pt}}}%
  \par\medskip}

% v2: 保护题干与开头，长题可跨页，续页补上下文。
\newcommand{\qbreak}{\needspace{6\baselineskip}}
% 独立练习：ID、来源、题干、考点、答案标签（通常 ans:ID）。
\newcommand{\practice}[5]{%
  \qbreak\subhdr{练习 #1}\qstem{（#2）#3}%
  \textbf{考点}：#4\par
  详解见附录第~\pageref{#5}~页。\par\medskip}
% 附录开头：ID、来源、条件摘要；自动产生 ans:ID 标签。
\newcommand{\answerstart}[3]{%
  \qbreak\subhdr{#1\quad 详解}\label{ans:#1}\qstem{（#2）#3}}
\newcommand{\answercontinue}[2]{%
  \needspace{5\baselineskip}\subhdr{#1\quad 续解}%
  {\small\color{inkgray}本页接续：#2}\par\medskip}
% 延迟提供默认值，让旧讲义可在导言区保留自己的 \newcommand{\step}。
\AtBeginDocument{\providecommand{\step}[1]{\needspace{3\baselineskip}\par\medskip\textbf{#1}\par}}

% 前提旁注：只放在「需要理解的那一步」后面，给陈述 + 本题代入，不重证。
\newcommand{\aside}[1]{%
  \par\vspace{1pt}\noindent
  \hangindent=1.6em\hangafter=1
  {\small\color{inkgray}%
    \hspace*{0.6em}{\color{petal}$\triangleright$}\hspace{2pt}#1}%
  \par\vspace{2pt}}

%---------------------------------------------------------------- 表格
% 本机 MiKTeX 的 colortbl 与 array 版本不匹配（\rowcolor 会报错），
% 故表格统一用「上粗紫线 + 表头 + 中细线」的三线式，不依赖 colortbl。
\newcommand{\tabtop}{\par\vspace{3pt}{\color{violet}\hrule height 1.2pt}\vspace{1pt}}
\newcommand{\tabmid}{\par\vspace{1pt}{\color{petal}\hrule height 0.5pt}\vspace{1pt}}
\newcommand{\tabbot}{\par\vspace{3pt}{\color{rulepurple}\hrule height 0.5pt}\vspace{8pt}}
\newcommand{\thead}[1]{{\bfseries\color{violetdeep}#1}}

%---------------------------------------------------------------- 页眉页脚
\pagestyle{fancy}
\fancyhf{}
\renewcommand{\headrulewidth}{0.4pt}
\renewcommand{\headrule}{\hbox to\headwidth{\color{rulepurple}\leaders\hrule height \headrulewidth\hfill}}
\fancyhead[L]{\small\color{inkgray}\songti \doctag}
\fancyhead[R]{\small\color{inkgray}\songti \coursetag}
\fancyfoot[C]{\small\color{inkgray}\orn\enspace{\color{violet}\thepage}\enspace\orn}

% 每篇讲义在导言区设这两个值
\newcommand{\doctag}{}
\newcommand{\coursetag}{}
\newcommand{\settags}[2]{\renewcommand{\doctag}{#1}\renewcommand{\coursetag}{#2}}

%---------------------------------------------------------------- 数学记号
\newcommand{\vect}[1]{\boldsymbol{#1}}
\newcommand{\dd}{\,\mathrm{d}}
\newcommand{\ddx}{\frac{\mathrm{d}}{\mathrm{d}x}}
\addtolength{\jot}{2pt}



\setlength{\headheight}{22pt}
\linespread{1.22}
\raggedbottom
\settags{专题 · 张量与并矢}{理论力学}

% 记号：矢量粗斜体，二阶张量粗正体
\newcommand{\ten}[1]{\mathbf{#1}}
\newcommand{\unitd}{\hat{\ten{1}}}
\newcommand{\step}[1]{\par\smallskip\noindent{\color{violetdeep}\bfseries #1}\hspace{0.35em}}

\let\orighdr\hdr
\renewcommand{\hdr}[2]{\par\needspace{5\baselineskip}\orighdr{#1}{#2}}
\let\origsubhdr\subhdr
\renewcommand{\subhdr}[1]{\par\needspace{3\baselineskip}\origsubhdr{#1}}

\setlength{\abovedisplayskip}{6pt plus 2pt minus 2pt}
\setlength{\belowdisplayskip}{6pt plus 2pt minus 2pt}

\begin{document}

\doctitle{从张量到并矢，再到理论力学}
{一条把“九个分量”讲成“一台机器”的路线}
{自编专题讲义\quad 2026-11-03\quad 接 Topic 8（刚体运动学）p8--p12}

\callout{\textbf{这份材料为谁写。}为“会算矩阵乘法，但二次型、迹、特征值都很生疏，张量完全没学过”的读者写。
它不讲三段彼此无关的知识，而是讲\textbf{一条路}：一个物理关系为什么必须是线性的（§1）$\to$ 线性映射的账本里有哪些换基也不变的数（§2）$\to$ 这些数怎么决定一个曲面（§3）$\to$ 同样的东西换个说法就叫张量（§4）$\to$ 再换个写法就叫并矢（§5）$\to$ 落到惯量张量、欧拉方程、应力（§7--§9）。
每一节只依赖它前面已经讲过的东西；凡是新概念，都配一个能口算的数字例子。}

\callout{\textbf{取材声明（哪些是课件的，哪些是补的）。}
课件 \texttt{TM2026-Topic8-P1.pdf} 直接给出的是：p8 的正交基与求和约定、$\varepsilon_{ijk}$ 与对偶张量 $\vect a^{\times}$；p9 的张量积（课件原词就叫 Dyad）、左右作用、转置、\textbf{把张量当线性算子 $b=Ta$}、\textbf{把张量当二次型 $E=aTa$}、单位张量 $\unitd$；p10--p11 的旋转算子 $S=\sum_k\vect n'_k\otimes\vect n_k$ 及其性质；p12 由 $S$ 导出的角速度。
\textbf{课件没有讲、由本讲义补的是}：为什么线性映射要用九个分量（§1）、迹与行列式为什么换基不变（§2）、二次型的等值面为什么是椭球、谱定理的完整陈述与证明思路（§3）、张量的变换律与商法则（§4）、并矢与矩阵/指标三套记号的对照（§6）。
凡补的内容都给出完整推导，不引“由某某定理”了事。}

\textbf{与旧版的关系。}本讲义取代 2026-11-03 的《专题·并矢：从张量积到惯量张量》。那一份从“为什么需要并矢”直接起步；它的并矢代数与三处落点在这里全部保留，但前面补了 §1--§4 的张量层，并把并置写法 $\vect a\vect b$ 收作正式别名（§5.1）。两处若有个别出入，以本讲义为准。

\hdr{§0}{路线图：三个坎，各自在哪节翻过去}

\begin{center}
\tabtop
\begin{tabular}{@{}p{0.13\linewidth}p{0.36\linewidth}p{0.42\linewidth}@{}}
\thead{坎} & \thead{你大概会卡在哪} & \thead{本讲义在哪里解决} \\
\multicolumn{3}{@{}l@{}}{\color{petal}\rule{\linewidth}{0.5pt}}\\
第一个 & “张量”听起来像新数学，其实是线性映射 & §1 把映射与矩阵分开；§4 说明张量就是“不依赖坐标的映射” \\
第二个 & 二次型、迹、特征值没见过，看不懂主轴 & §2--§3 从零补：迹为什么不变、二次型的等值面是椭球、谱定理 \\
第三个 & 课堂上老师写 $ab$ 或 $\otimes$，不知道怎么算 & §5 给出作用规则与三套记号的字典；§10 三道例题跑一遍 \\
\end{tabular}
\tabbot
\end{center}

\textbf{读法。}如果你时间紧，最短路径是 §1.2 $\to$ §2.1 $\to$ §3.3 $\to$ §4.3 $\to$ §5.1 $\to$ §7.1。这条路上不跳过任何“后面才用”的定义，但会跳过所有例子。
反过来，如果你手上有课件 p9 那页想对上号，直接翻 §6 的三记号字典，再回 §5 补作用规则。

\warnbox{\textbf{全篇只有一个主线，先记住它。}
\textbf{同一个物理对象，换坐标系时分量会变，但对象本身不变。}
张量是“对象”，矩阵是“这个对象在一组基下的账本”，并矢是“这个对象的另一种写法”。
后面每一节都在这句话上做文章：§1 说清“账本”是什么，§2 找出账本里不随基变的数，§4 用“怎么变”来定义张量，§5 用“怎么写”来避开矩阵。}

\hdr{§1}{先补一件事：矩阵到底是什么}

\subhdr{1.1\quad 线性映射：矢量进、矢量出}

理论力学里到处是“给一个矢量，吐一个矢量”的关系：

\begin{center}
\tabtop
\begin{tabular}{@{}lll@{}}
\thead{关系} & \thead{进} & \thead{出} \\
\multicolumn{3}{@{}l@{}}{\color{petal}\rule{\linewidth}{0.5pt}}\\
功 $W=\vect F\cdot\vect r$ & 矢量 $\vect r$ & \textbf{数}（不是矢量） \\
力矩 $\vect M=\vect r\times\vect F$ & 两个矢量 & 矢量 \\
角动量 $\vect L=\ten J\cdot\vect\omega$ & 角速度 $\vect\omega$ & 矢量 \\
\end{tabular}
\tabbot
\end{center}

第一行吃矢量吐数，第二三行吃矢量吐矢量。\textbf{本讲义只关心后一类}，因为只有它们需要“九个分量”这种对象。

\begin{center}
\fcolorbox{violet}{bandviolet}{\parbox{0.94\linewidth}{\centering\vspace{3pt}
\small 定义：映射 $T$ 若满足 $T(\alpha\vect u+\beta\vect v)=\alpha T(\vect u)+\beta T(\vect v)$ 对任意 $\vect u,\vect v$ 与任意数 $\alpha,\beta$ 成立，就称 $T$ 是\textbf{线性映射}。\vspace{3pt}}}
\end{center}

为什么要这条性质？因为它省事到极点：\textbf{只要知道 $T$ 作用在三个基矢量上的结果，就知道了它对一切矢量的结果。} 理由是任意矢量都能写成 $\vect x=x_1\vect e_1+x_2\vect e_2+x_3\vect e_3$，于是

\[
T(\vect x)=T\Bigl(\sum_j x_j\vect e_j\Bigr)=\sum_j x_j\,T(\vect e_j).
\]

右边只用到三个东西 $T(\vect e_1),T(\vect e_2),T(\vect e_3)$。每个 $T(\vect e_j)$ 本身是一个矢量，需要三个数才能写出来，所以\textbf{一共九个数的信息量}。
这就是“二阶张量有九个分量”的全部来源——不是玄学，是“三个输入基矢量，每个吐三个分量”。

\subhdr{1.2\quad 矩阵是这台机器的账本}

把 $T(\vect e_j)$ 在基 $\{\vect e_i\}$ 下展开，系数记作 $T_{ij}$：

\[
T(\vect e_j)=\sum_{i=1}^{3}T_{ij}\,\vect e_i
\qquad\Longleftrightarrow\qquad
\boxed{\;(T\vect x)_i=\sum_{j=1}^{3}T_{ij}x_j\;}
\]

把九个 $T_{ij}$ 按“第 $i$ 行、第 $j$ 列”排成方阵，就是 $T$ 的矩阵。\textbf{请特别注意指标的位置：第二个指标 $j$ 是“输入基矢量”的编号，第一个 $i$ 是“输出分量”的编号。} 这个约定决定了后面所有点乘的方向，混了就全错。

\textbf{数字例子。}设 $T$ 把 $\vect e_1$ 变成 $2\vect e_1$，把 $\vect e_2$ 变成 $\vect e_1+\vect e_2$，把 $\vect e_3$ 变成 $3\vect e_3$：

\[
T(\vect e_1)=\begin{pmatrix}2\\0\\0\end{pmatrix},\quad
T(\vect e_2)=\begin{pmatrix}1\\1\\0\end{pmatrix},\quad
T(\vect e_3)=\begin{pmatrix}0\\0\\3\end{pmatrix}
\quad\Longrightarrow\quad
\ten T=\begin{pmatrix}2&1&0\\0&1&0\\0&0&3\end{pmatrix}.
\]

注意这三个列矢量是按列摆进去的。现在取 $\vect x=(1,2,3)$：

\[
\ten T\vect x=\begin{pmatrix}2&1&0\\0&1&0\\0&0&3\end{pmatrix}\begin{pmatrix}1\\2\\3\end{pmatrix}
=\begin{pmatrix}2\cdot1+1\cdot2\\0\cdot1+1\cdot2\\0+0+3\cdot3\end{pmatrix}=\begin{pmatrix}4\\2\\9\end{pmatrix}.
\]

口算核对：$T(\vect x)=\vect x_1T(\vect e_1)+\vect x_2T(\vect e_2)+\vect x_3T(\vect e_3)=1\cdot(2,0,0)+2\cdot(1,1,0)+3\cdot(0,0,3)=(4,2,9)$。两条路一致。

\subhdr{1.3\quad 换基：账本变，机器不变}

同一台机器，换一组基来记账，九个数字会变。这句话是本讲义的枢纽，值得用一个具体例子钉死。

\textbf{二维例子。}机器 $T$ 是“沿 $x$ 轴拉伸两倍”：$T(x,y)=(2x,y)$。在标准基 $\vect e_1,\vect e_2$ 下

\[
\ten T=\begin{pmatrix}2&0\\0&1\end{pmatrix}.
\]

现在改用转了 $45^\circ$ 的基 $\vect n_1=(\vect e_1+\vect e_2)/\sqrt2$，$\vect n_2=(-\vect e_1+\vect e_2)/\sqrt2$。
先算输入：$T\vect n_1=T\bigl(\tfrac{1}{\sqrt2},\tfrac{1}{\sqrt2}\bigr)=\bigl(\tfrac{2}{\sqrt2},\tfrac{1}{\sqrt2}\bigr)$。
再把它在这组新基下分解，用 $\vect n_i\cdot$ 投影取系数：

\[
T\vect n_1\cdot\vect n_1=\frac{2+1}{2}=\frac32,\qquad
T\vect n_1\cdot\vect n_2=\frac{-2+1}{2}=-\frac12 .
\]

同法 $T\vect n_2=T\bigl(-\tfrac{1}{\sqrt2},\tfrac{1}{\sqrt2}\bigr)=\bigl(-\tfrac{2}{\sqrt2},\tfrac{1}{\sqrt2}\bigr)$，

\[
T\vect n_2\cdot\vect n_1=\frac{-2+1}{2}=-\frac12,\qquad
T\vect n_2\cdot\vect n_2=\frac{2+1}{2}=\frac32 .
\]

于是同一台机器的账本变成

\[
\ten T'=\begin{pmatrix}3/2&-1/2\\-1/2&3/2\end{pmatrix}.
\]

两个账本长得完全不一样，可它们描述的是同一件事：沿 $\vect e_1$ 方向拉两倍、沿 $\vect e_2$ 方向不变。
\textbf{判断“是不是同一台机器”，不能看数字，要看数字怎么变。}这条换算规则是

\[
\boxed{\;T'_{ij}=\sum_{k,l}R_{ki}R_{lj}T_{kl}\;}
\qquad\text{其中 }R_{ki}=\vect n_i\cdot\vect e_k .
\]

$R$ 的每一列是新基矢量在旧基下的分量；正交基之间 $R^{-1}=R^{\mathsf T}$（转置），所以上式也常写成 $\ten T'=\ten R^{\mathsf T}\ten T\ten R$。

\aside{上面那个例子里 $R=\frac{1}{\sqrt2}\begin{pmatrix}1&-1\\1&1\end{pmatrix}$。代入 $\ten R^{\mathsf T}\ten T\ten R$ 得 $\begin{pmatrix}3/2&-1/2\\-1/2&3/2\end{pmatrix}$，与逐项算的一致。这里不重证，只用来核对上面的两次投影没算错。}

\subhdr{1.4\quad 账本里有没有不随基变的数}

既然换个基数字全变，一个自然的问题是：\textbf{账本里有没有哪个组合换基后不变？}
有，而且恰好是需要的那几个。这就是迹和行列式，也是 §3 讲主轴的入口。

先记住这个念头的形状：\textbf{不变量 = 与坐标系无关的物理量的载体。} 转动惯量的三个主值、应力的三个主应力，都是这类数。

\hdr{§2}{迹与行列式：两个换基也不变的数}

\subhdr{2.1\quad 迹：对角元之和，为什么换基不变}

\begin{center}
\fcolorbox{violet}{bandviolet}{\parbox{0.94\linewidth}{\centering\vspace{3pt}
\small 定义：方阵 $\ten A$ 的\textbf{迹}是对角元之和，$\operatorname{tr}\ten A=\sum_i A_{ii}$。\vspace{3pt}}}
\end{center}

\textbf{数字例子。}上面 §1.3 的两个账本：

\[
\operatorname{tr}\begin{pmatrix}2&0\\0&1\end{pmatrix}=2+1=3,
\qquad
\operatorname{tr}\begin{pmatrix}3/2&-1/2\\-1/2&3/2\end{pmatrix}=\frac32+\frac32=3 .
\]

一样。这不是巧合。

\textbf{为什么不变：一句话加一步代数。}矩阵乘法有一条“循环性”：对任意两个方阵，

\[
\operatorname{tr}(\ten A\ten B)=\sum_i\sum_j A_{ij}B_{ji}=\sum_j\sum_i B_{ji}A_{ij}=\operatorname{tr}(\ten B\ten A).
\]

于是，令 $\ten B=\ten R$、$\ten A=\ten R^{\mathsf T}\ten T$（这样 $\ten R^{\mathsf T}\ten T\ten R=\ten A\ten B$）：

\[
\operatorname{tr}(\ten R^{\mathsf T}\ten T\ten R)=\operatorname{tr}\bigl((\ten R^{\mathsf T}\ten T)\ten R\bigr)
=\operatorname{tr}\bigl(\ten R(\ten R^{\mathsf T}\ten T)\bigr)
=\operatorname{tr}\bigl((\ten R\ten R^{\mathsf T})\ten T\bigr)
=\operatorname{tr}\ten T .
\]

最后一步用了 $\ten R\ten R^{\mathsf T}=\unitd$（正交基的性质，见 §1.3）。
\textbf{一行读完：迹不变，因为迹允许把因子转圈，而正交矩阵转一圈就消掉了。}

\subhdr{2.2\quad 行列式：体积的缩放倍数}

行列式在这里只需要一个图像，不需要会算。\textbf{$\det\ten T$ 是这台机器对体积的缩放倍数。}
把单位立方体喂给 $T$，出来一个平行六面体，它的体积就是 $|\det T|$；符号表示“是否翻转了手性”。

这个图像立刻给出两个结论：
\begin{itemize}
\item 换坐标系不会改变“体积被放大几倍”这件事，所以 $\det$ 也是不变量。代数上同样一句话：

\[
\det(\ten R^{\mathsf T}\ten T\ten R)=\det(\ten R^{\mathsf T})\det\ten T\det(\ten R)=\det\ten T,
\]

因为正交矩阵满足 $\det\ten R=\pm1$，平方为 $1$。
\item $\det\ten T=0$ 意味着\textbf{体积被压成零}，即这台机器把三维压成平面或直线，不可逆。后面“某个主值为零”就是这种退化。
\end{itemize}

\textbf{数字例子。}§1.3 里 $\ten T=\operatorname{diag}(2,1)$，$\det\ten T=2$：面积放大两倍（横向拉两倍、纵向不变）。
换基后

\[
\ten T'=\begin{pmatrix}3/2&-1/2\\-1/2&3/2\end{pmatrix},
\qquad
\det\ten T'=\frac94-\frac14=2 ,
\]

一样。

\subhdr{2.3\quad 特征多项式：把不变量一网打尽}

把“$\det$ 与 $\operatorname{tr}$ 都不变”这个念头推到底，就是问：$\det(\ten T-\lambda\unitd)$ 作为 $\lambda$ 的多项式是什么？它的系数是不变量，根就是特征值。

二维先看清楚。$\ten T=\begin{pmatrix}a&b\\c&d\end{pmatrix}$，

\[
\det(\ten T-\lambda\unitd)=\det\begin{pmatrix}a-\lambda&b\\c&d-\lambda\end{pmatrix}
=(a-\lambda)(d-\lambda)-bc=\lambda^2-(a+d)\lambda+(ad-bc).
\]

所以二维的式子只有两个系数，分别就是 $\operatorname{tr}$ 与 $\det$：

\[
\lambda^2-(\operatorname{tr}\ten T)\lambda+\det\ten T=0 .
\]

三维照做，得到一个三次多项式：

\[
\det(\ten T-\lambda\unitd)=-\lambda^3+I_1\lambda^2-I_2\lambda+I_3 ,
\]
\[
I_1=\operatorname{tr}\ten T,\qquad
I_3=\det\ten T,\qquad
I_2=\tfrac12\bigl[(\operatorname{tr}\ten T)^2-\operatorname{tr}\ten T^2\bigr].
\]

三个系数 $I_1,I_2,I_3$ 都由 $\operatorname{tr}$ 与 $\det$ 拼成，所以\textbf{三个都不随基变}。
这三个数就是 §3.5 要用的“三个不变量”，现在先知道它们从哪来就够了。

\warnbox{\textbf{为什么要有 $\ten T^2$。}初学者常问 $I_2$ 里为什么要平方。看二维就明白了：$\operatorname{tr}\ten T^2=a^2+2bc+d^2$，而 $(\operatorname{tr}\ten T)^2=a^2+2ad+d^2$，两者相减再除二得 $ad-bc=\det\ten T$，正是二维的“第二个不变量”。
所以 $I_2$ 不是新东西，是\textbf{二维里 $\det$ 的位置在三维里的替身}——三维有三个独立的不变量，需要三个组合。}

\hdr{§3}{二次型：课件 p9 那句 $E=aTa$ 到底是什么}

\subhdr{3.1\quad 从课件那一行说起}

课件 p9 最后一行写的是

\[
\text{Tensor as a quadratic form:}\qquad E=\vect a\ten T\vect a=T_{ik}a_ia_k .
\]

这句话的读法只有一层：\textbf{给它一个矢量，吐出一个数，而且这个数对 $\vect a$ 是二次的。}
“二次的”意思是：把 $\vect a$ 换成 $2\vect a$，$E$ 变成 $4E$；换成 $-\vect a$，$E$ 不变。所以它刻画的是“方向 + 大小”，不是“矢量”。

三维写开就是

\[
E=\sum_{i=1}^{3}\sum_{k=1}^{3}T_{ik}a_ia_k
=T_{11}a_1^2+T_{22}a_2^2+T_{33}a_3^2+2T_{12}a_1a_2+2T_{13}a_1a_3+2T_{23}a_2a_3 .
\]

注意交叉项前面那个 $2$：它是 $T_{12}a_1a_2$ 与 $T_{21}a_2a_1$ 两项合并的结果。这个 $2$ 后面会反复出现，别丢。

\textbf{物理上它长什么样。}刚体的转动动能

\[
T=\tfrac12\vect\omega\cdot\ten J\cdot\vect\omega=\tfrac12J_{ik}\omega_i\omega_k
\]

就是二次型（差一个 $\tfrac12$）。这句话解释了为什么“惯量”只能用张量描述：能量是 $\vect\omega$ 的二次函数，而二次函数的系数天然是一个对称的九元素表。

\subhdr{3.2\quad 只有对称部分起作用}

\textbf{一条必须先钉死的事实}：二次型对 $\ten T$ 的\textbf{反对称部分完全无感}。也就是说，算 $E$ 时可以把 $\ten T$ 换成它自己的对称部分，结果一字不差。

理由是纯粹的指标游戏。交叉项 $T_{ik}a_ia_k$ 里，把 $i$ 与 $k$ 对换名字（只是换字母，值不变）：

\[
\sum_{i,k}T_{ik}a_ia_k=\sum_{i,k}T_{ki}a_ka_i=\sum_{i,k}T_{ki}a_ia_k .
\]

所以它等于 $\sum_{i,k}\tfrac12(T_{ik}+T_{ki})a_ia_k$，而 $\tfrac12(\ten T+\ten T^{\mathsf T})$ 正是对称部分。

\textbf{数字例子。}取 $\ten T=\begin{pmatrix}1&2\\0&1\end{pmatrix}$（不对称）。算 $E$：

\[
E=\begin{pmatrix}x&y\end{pmatrix}\begin{pmatrix}1&2\\0&1\end{pmatrix}\begin{pmatrix}x\\y\end{pmatrix}
=x^2+2xy+y^2=(x+y)^2 .
\]

换成对称部分 $\tfrac12(\ten T+\ten T^{\mathsf T})=\begin{pmatrix}1&1\\1&1\end{pmatrix}$：

\[
E=\begin{pmatrix}x&y\end{pmatrix}\begin{pmatrix}1&1\\1&1\end{pmatrix}\begin{pmatrix}x\\y\end{pmatrix}=x^2+xy+yx+y^2=x^2+2xy+y^2 .
\]

一样。\textbf{所以从这一节起，凡讲二次型都默认 $\ten T$ 对称。}力学里的 $\ten J$ 与 $\ten\sigma$ 本来就对称（§7.2、§9.2 会各给一次理由），不会损失任何东西。

\subhdr{3.3\quad 二次型的等值面是椭球：主轴就是它的轴}

令 $E=1$，得到一个曲面。这个曲面的形状决定了一切。

\textbf{没有交叉项的情形。}$\ten T=\begin{pmatrix}1&0\\0&4\end{pmatrix}$，则 $E=x^2+4y^2=1$，这是椭圆

\[
\frac{x^2}{1^2}+\frac{y^2}{(1/2)^2}=1,
\]

半轴分别是 $1$ 与 $1/2$。注意半轴是 $1/\sqrt{\lambda}$，$\lambda$ 是 $\ten T$ 的对角元（也是它的特征值）：$\lambda=1$ 给半轴 $1$，$\lambda=4$ 给半轴 $1/2$。\textbf{特征值越大，那个方向越“扁”。}

\textbf{有交叉项的情形。}$\ten T=\begin{pmatrix}2&1\\1&2\end{pmatrix}$，则 $E=2x^2+2xy+2y^2=1$。这个椭圆是斜的。把坐标转 $45^\circ$：

\[
x=\frac{u+v}{\sqrt2},\qquad y=\frac{-u+v}{\sqrt2}
\quad\Longrightarrow\quad
x^2+y^2=u^2+v^2,\qquad xy=\frac{u^2-v^2}{2},
\]

代入：

\[
E=2(u^2+v^2)+2\cdot\frac{u^2-v^2}{2}=3u^2+v^2 .
\]

\textbf{交叉项消失了}，曲面变成半轴 $1/\sqrt3$ 与 $1$ 的正椭圆。让交叉项消失的那组坐标轴 $u,v$，就是\textbf{主轴}。

这不是巧合，而是全部内容：\textbf{对任何对称 $\ten T$，总存在一组正交坐标轴，使二次型没有交叉项；在那一组轴上，系数就是特征值。}
这句话的严格版本就是 §3.5 的谱定理。

\aside{$45^\circ$ 这个角不是猜的：$\ten T$ 的两个特征向量是 $(1,1)/\sqrt2$ 与 $(-1,1)/\sqrt2$，与 $x$ 轴成 $45^\circ$ 与 $135^\circ$。凡是先算出特征向量再取转角，都会得到这个数。}

\subhdr{3.4\quad 特征值与特征向量：主轴方向的定义}

\begin{center}
\fcolorbox{violet}{bandviolet}{\parbox{0.94\linewidth}{\centering\vspace{3pt}
\small 若存在非零矢量 $\vect n$ 与数 $\lambda$ 使 $\ten T\vect n=\lambda\vect n$，则称 $\lambda$ 为\textbf{特征值}、$\vect n$ 为对应的\textbf{特征向量}。\vspace{3pt}}}
\end{center}

几何含义：\textbf{$\ten T$ 作用在 $\vect n$ 上不改变方向，只改变长度。} 这就是“主轴”三个字的全部意思——沿这根轴，机器只做伸缩。

怎么求？把 $\ten T\vect n=\lambda\vect n$ 移项成 $(\ten T-\lambda\unitd)\vect n=\vect 0$。这是齐次线性方程组，要它有\textbf{非零}解，系数矩阵必须不可逆（否则唯一解是 $\vect n=\vect0$），也就是

\[
\det(\ten T-\lambda\unitd)=0 .
\]

这正是 §2.3 那个多项式等于零。所以求特征值就是解一个三次方程——\textbf{而它的三个系数正是三个不变量}。

\textbf{为什么主轴方向的法向与径向平行。}椭圆 $E=1$ 上一点 $\vect x$，其法向是 $\nabla E=2\ten T\vect x$。若 $\vect x$ 落在主轴方向上，由对称性这根轴上的法向必须与 $\vect x$ 平行（否则椭圆在这点“歪”了，就不是轴），于是 $\ten T\vect x\parallel\vect x$——这就是特征值方程。
反过来说：\textbf{“主轴方向”与“特征向量”是同一件事的两种说法}，一个来自几何，一个来自代数。

\subhdr{3.5\quad 谱定理：完整陈述与证明思路}

\begin{center}
\fcolorbox{violet}{bandviolet}{\parbox{0.94\linewidth}{\vspace{3pt}
\textbf{谱定理（实对称矩阵，三维）。}设 $\ten T$ 是 $3\times3$ 实矩阵且 $\ten T=\ten T^{\mathsf T}$。则
\begin{enumerate}[label=(\roman*),leftmargin=1.8em,itemsep=2pt]
\item $\ten T$ 的三个特征值 $\lambda_1,\lambda_2,\lambda_3$ 全是实数；
\item 属于不同特征值的特征向量互相正交；
\item 存在三个互相正交的单位矢量 $\vect n_1,\vect n_2,\vect n_3$（三个主轴），使
\[
\ten T=\lambda_1\vect n_1\vect n_1^{\mathsf T}+\lambda_2\vect n_2\vect n_2^{\mathsf T}+\lambda_3\vect n_3\vect n_3^{\mathsf T},
\]
即：在以 $\{\vect n_k\}$ 为基的坐标系里，$\ten T$ 的账本就是 $\operatorname{diag}(\lambda_1,\lambda_2,\lambda_3)$。
\end{enumerate}
\vspace{3pt}}}
\end{center}

\textbf{(ii) 的证明只有两行。}设 $\ten T\vect u=\lambda\vect u$，$\ten T\vect v=\mu\vect v$，$\lambda\ne\mu$。则

\[
\lambda(\vect u\cdot\vect v)=(\ten T\vect u)\cdot\vect v=\vect u\cdot(\ten T^{\mathsf T}\vect v)=\vect u\cdot(\ten T\vect v)=\mu(\vect u\cdot\vect v).
\]

第二个等号用了“转置的判别式”$(\ten T\vect u)\cdot\vect v=\vect u\cdot(\ten T^{\mathsf T}\vect v)$（§5.5 会正式给这条），第三个用了对称 $\ten T^{\mathsf T}=\ten T$。于是 $(\lambda-\mu)(\vect u\cdot\vect v)=0$，而 $\lambda\ne\mu$，只能 $\vect u\cdot\vect v=0$。

\textbf{(i)+(iii) 的证明思路（拉格朗日乘子，值得跟一遍）。}
考虑在单位球面 $\vect x\cdot\vect x=1$ 上求二次型 $f(\vect x)=\vect x\cdot\ten T\vect x$ 的最大值。
球面是闭且有界的，$f$ 连续，所以最大值一定取到，记在 $\vect n_1$ 处。用乘子法：要 $\nabla f=\mu\nabla g$，其中 $g=\vect x\cdot\vect x-1$。而

\[
\nabla f=2\ten T\vect x,\qquad \nabla g=2\vect x
\quad\Longrightarrow\quad
\ten T\vect n_1=\mu\vect n_1 .
\]

所以极值点必是特征向量，乘子 $\mu$ 就是特征值；并且在该点 $f=\vect n_1\cdot\ten T\vect n_1=\mu\,|\vect n_1|^2=\mu$。
\textbf{一句话：单位球面上二次型的最大值，就是最大的特征值。}
接着在“与 $\vect n_1$ 正交的那个平面”里重复同样的论证（维数降一），得到 $\vect n_2$；再降一次得到 $\vect n_3$。
每一步都只在实空间里取极值，所以得到的 $\lambda$ 必是实数——这就是 (i)。

\aside{这条证明同时给出了一个以后天天用的结论：$\lambda_{\min}\le \vect x\cdot\ten T\vect x\le\lambda_{\max}$ 对一切单位矢量成立（瑞利商）。§7 用它来说明“沿哪个方向转，转动惯量最大”。}

\subhdr{3.6\quad 三个不变量：验算的最好工具}

由 §2.3，三个系数是

\[
I_1=\operatorname{tr}\ten T=\lambda_1+\lambda_2+\lambda_3,
\qquad
I_3=\det\ten T=\lambda_1\lambda_2\lambda_3,
\]
\[
I_2=\tfrac12\bigl[(\operatorname{tr}\ten T)^2-\operatorname{tr}\ten T^2\bigr]
=\lambda_1\lambda_2+\lambda_2\lambda_3+\lambda_3\lambda_1 .
\]

右边那两个等号（用 $\lambda$ 表示）不需要另证：在主轴系里 $\ten T=\operatorname{diag}(\lambda_1,\lambda_2,\lambda_3)$，而 $\operatorname{tr}$ 与 $\det$ 都不随基变（§2），直接代进去就是。

\textbf{用法：做完一道特征值题，用 $I_1$ 与 $I_3$ 各查一遍。}以 §3.3 的 $\ten T=\begin{pmatrix}2&1\\1&2\end{pmatrix}$ 为例：$\operatorname{tr}=4$、$\det=3$，而算出的特征值是 $3$ 与 $1$，$3+1=4$、$3\cdot1=3$，对上。

三维情形同样：§10 例 1 里 $\operatorname{tr}=10$、$\det=30$，特征值 $5,3,2$ 满足 $5+3+2=10$、$5\cdot3\cdot2=30$。

\warnbox{\textbf{一个容易误解的地方。}“$I_1,I_2,I_3$ 不变”说的是\textbf{换基时}不变，不是“$\ten T$ 不变”。换基时九个分量全变，只有这三个组合不变。
另外，特征值不变但\textbf{特征向量会变}——它随基一起转。所以“主轴方向”这句话必须连着坐标系说，否则没有意义。}

\hdr{§4}{张量：从多重线性到变换律}

\subhdr{4.1\quad 阶：拿几个矢量喂它}

“张量”这个词让人紧张，但它的定义方式很朴素：\textbf{按它需要几个矢量才能给出一个数，来分阶。}

\begin{center}
\tabtop
\begin{tabular}{@{}p{0.12\linewidth}p{0.40\linewidth}p{0.40\linewidth}@{}}
\thead{阶} & \thead{吃什么、吐什么} & \thead{例子} \\
\multicolumn{3}{@{}l@{}}{\color{petal}\rule{\linewidth}{0.5pt}}\\
0 阶 & 不吃，直接是一个数 & 温度、质量、动能 \\
1 阶 & 吃一个矢量，吐一个数；且对输入线性 & 力在方向 $\vect u$ 上的投影 $\vect F\cdot\vect u$ \\
2 阶 & 吃两个矢量，吐一个数；且对每个输入都线性 & 应力：$\vect u\cdot\ten\sigma\cdot\vect v$；惯量：$\vect u\cdot\ten J\cdot\vect v$ \\
\end{tabular}
\tabbot
\end{center}

“对每个输入都线性”叫\textbf{多重线性}。对二阶就是要

\[
\ten T(\alpha\vect u_1+\beta\vect u_2,\;\vect v)=\alpha\ten T(\vect u_1,\vect v)+\beta\ten T(\vect u_2,\vect v),
\]

对第二个位置同样要求（两个位置各要求一次，合起来叫“双线性”）。

\textbf{为什么二阶张量有两个身份。}双线性的 $\ten T(\vect u,\vect v)$ 可以“先喂一个”：

\[
\vect u\;\longmapsto\;\bigl[\vect v\mapsto\ten T(\vect u,\vect v)\bigr] .
\]

右边是一个“吃 $\vect v$ 吐数”的线性函数，在欧氏空间里它就是某个矢量（用内积可以一一对应，见 §4.2）。所以

\[
\boxed{\;\text{吃两个矢量吐一个数}\quad\Longleftrightarrow\quad\text{吃一个矢量吐一个矢量}\;}
\]

这两个说法是同一件事。课件 p9 写的是右边这个（$b=Ta$）；本讲义 §3 用的是左边这个（$E=\vect a\ten T\vect a$）。\textbf{看起来是两套东西，其实是同一个对象的两个面。}

\subhdr{4.2\quad 张量积：并矢的出生地}

现在造一个具体的双线性对象。给两个矢量 $\vect a,\vect b$，规定

\[
(\vect a\otimes\vect b)(\vect u,\vect v)\stackrel{\text{def}}{=}(\vect a\cdot\vect u)(\vect b\cdot\vect v).
\]

右边是两个数的乘积，所以确实吐出一个数；对 $\vect u$ 线性、对 $\vect v$ 线性，都满足。
这个对象记作 $\vect a\otimes\vect b$，叫 $\vect a$ 与 $\vect b$ 的\textbf{张量积}。\textbf{课件 p9 给它起的名字是 Dyad，中文就叫并矢。}

它跟“喂一个矢量”的那个身份怎么对上？利用上面的等价，把第二个位置定死为 $\vect v$：

\[
(\vect a\otimes\vect b)\cdot\vect v\;\stackrel{\text{def}}{=}\;\vect a\,(\vect b\cdot\vect v).
\]

这就是全部的作用规则，一句话：\textbf{先把 $\vect v$ 投影到 $\vect b$ 上，得到一个数，再拿这个数去放大 $\vect a$。}
（另一边 $\vect v\cdot(\vect a\otimes\vect b)=(\vect v\cdot\vect a)\vect b$ 是另一条规则，二者一般不同，§5.2 专门讲。）

\textbf{九个基并矢。}取正交基 $\{\vect e_i\}$，把 $i,j$ 各取 $1,2,3$，得到九个对象 $\vect e_i\otimes\vect e_j$。它们是线性无关的，所以任何二阶张量都能唯一展开成

\[
\ten T=\sum_{i=1}^{3}\sum_{j=1}^{3}T_{ij}\,\vect e_i\otimes\vect e_j .
\]

\textbf{这就是“九个分量”的来历}：九个基并矢的系数，一个不多一个不少。
对照 §1.1 的论证——三个输入基矢量、每个吐三个分量——是同一件事的两种数法。

\subhdr{4.3\quad 指标写法与求和约定}

课件 p8 已经给了求和约定：\textbf{同一项里重复出现两次的指标自动求和，不写 $\sum$。}本讲义正文照旧写求和号（免得看错缩并了哪个指标），只在 §6 的字典与 §7 之后的推导里用约定。

有了约定，上面那些式子变得很短：

\[
\ten T=T_{ij}\,\vect e_i\otimes\vect e_j,\qquad
(\ten T\vect x)_i=T_{ij}x_j,\qquad
E=T_{ij}a_ia_j .
\]

$\delta_{ij}$（$i=j$ 时为 $1$、否则为 $0$）叫\textbf{克罗内克符号}，它对应单位张量：$\delta_{ij}\vect e_i\otimes\vect e_j=\unitd$。

\textbf{张量积的分量就是分量的乘积。}这条要记牢：

\[
(\vect a\otimes\vect b)_{ij}=a_ib_j
\qquad\Longrightarrow\qquad
\vect a\otimes\vect b=
\begin{pmatrix}
a_1b_1&a_1b_2&a_1b_3\\
a_2b_1&a_2b_2&a_2b_3\\
a_3b_1&a_3b_2&a_3b_3
\end{pmatrix}.
\]

矩阵的样子是“\textbf{列乘行}”，也叫外积。取 $\vect a=(1,2,0)$、$\vect b=(0,1,3)$，则

\[
\vect a\otimes\vect b=\begin{pmatrix}0&1&3\\0&2&6\\0&0&0\end{pmatrix}.
\]

\textbf{注意它不对称}：$(1,2)$ 元是 $1$，$(2,1)$ 元是 $0$。这就是“顺序有意义”的直接证据。

\subhdr{4.4\quad 变换律：张量与“九个数”的分界线}

现在回答一个关键问题：\textbf{随便给我九个排成方阵的数，它是不是张量？}

答案：不一定。张量的定义不在“有九个分量”，而在“换基时这九个分量怎么变”。规则就是 §1.3 那条：

\[
\boxed{\;T'_{ij}=R_{ki}R_{lj}T_{kl}\;}
\qquad\text{（$R_{ki}=\vect n_i\cdot\vect e_k$ 是新基在旧基下的分量）}
\]

凡按这条规则变的九元数组，才是张量；不按这条变的，只是九个数字。

\textbf{一个真的张量：$\delta_{ij}$。}把 $T_{kl}=\delta_{kl}$ 代进去：

\[
\delta'_{ij}=R_{ki}R_{lj}\delta_{kl}=R_{ki}R_{kj}=(\ten R\ten R^{\mathsf T})_{ij}=\delta_{ij}.
\]

用到了正交性 $\ten R\ten R^{\mathsf T}=\unitd$。结果说明：\textbf{单位张量在任何正交基下的分量都是 $\delta_{ij}$}。这正是“单位矩阵换基后还是单位矩阵”的深层原因，也是课件 p9 那句 $\unitd=\vect n_k\otimes\vect n_k$ 对任意正交基都成立的理由。

\textbf{一个不是张量的东西：位置矢量。}三个数 $(x_1,x_2,x_3)$ 在\textbf{旋转}下确实按矢量规则变（所以转动问题里可以当矢量用），但一旦允许\textbf{平移}坐标系，它会多出一个常数项，不再是齐次的。
所以“是不是张量”这句话必须连着\textbf{允许哪些变换}说。本讲义（以及整门理论力学）只允许旋转与反射，即\textbf{正交变换}，在这个群里位置矢量是矢量。

\warnbox{\textbf{这条区分有什么用。}它把“张量”从名词变成了\textbf{可检验的命题}。
以后看到任何一组九个量，只要你能证明它在换基时按上式变，它就是张量，于是 §3 的谱定理、不变量、主轴全套都可以直接搬过去用；证不出来，就别乱用。}

\subhdr{4.5\quad 商法则：从物理定律反认张量}

实践中更常用的是反过来：手上有一个物理关系，想证明里面那个九元数组是张量。

\begin{center}
\fcolorbox{violet}{bandviolet}{\parbox{0.94\linewidth}{\vspace{3pt}
\textbf{商法则（三维，实用版）。}若对\textbf{任意}矢量 $\vect v$ 都有 $u_i=T_{ij}v_j$，且已知 $\vect u$、$\vect v$ 都是矢量（即在旋转下按矢量规则变），则 $T_{ij}$ 必是二阶张量。
\vspace{3pt}}}
\end{center}

道理很直接：在带撇的坐标系里同一条定律必须同样成立，$u'_i=T'_{ij}v'_j$；把 $u'=Ru$、$v'=Rv$ 代进去，与 $u=TRv$ 比较，就逼出 $T'=RTR^{\mathsf T}$，正是 §4.4 的变换律。

\textbf{它马上要用在哪。}§9 里“面力是法向的函数”写作 $\vect t(\vect n)=\ten\sigma\cdot\vect n$。
$\vect t$ 是矢量、$\vect n$ 是矢量，所以由商法则，$\ten\sigma$ 一定是二阶张量——\textbf{不是“我们假设它是张量”，而是被这条定律逼出来的。}这是全校理论力学里最漂亮的一步推理。

\subhdr{4.6\quad 预告：上下标是从哪来的}

你以后会看到 $T^{i}{}_{j}$、$T_{ij}$、$T^{ij}$ 三种写法并存。这里给半页说明，本讲义正文不用。

在\textbf{正交}基里，三种写法数值完全相同，所以随便挑一种。但在\textbf{非正交}基（斜坐标、曲线坐标）里，矢量有两套分量：
把矢量投影到基矢量上得到一套（协变，下标记），把矢量沿基矢量分解得到另一套（逆变，上标记），二者由一个叫度规的 $g_{ij}$ 互相转换。
只有当 $g_{ij}=\delta_{ij}$（正交归一）时，两套才重合，上下标才可以不分。

\textbf{结论：}理论力学的刚体部分全程正交基，所以本讲义一律用下标，不会错。
等到连续介质力学用曲线坐标、或相对论用四维时空时，上下标就必须写清楚——那时再回来补这半页即可。

\hdr{§5}{并矢：Gibbs 的那套写法}

\textbf{这一节没有新数学，只换写法。}§4 已经说明二阶张量是双线性对象；现在把它写成不用矩阵的样子。

\subhdr{5.1\quad 并置：$ab$ 与 $\vect a\otimes\vect b$ 是同一个东西}

Gibbs 在 1880 年代引入这套记号时，把张量积\textbf{直接并排写}：

\[
\vect a\vect b\;\equiv\;\vect a\otimes\vect b,
\qquad\text{作用规则}\quad
(\vect a\vect b)\cdot\vect x=\vect a\,(\vect b\cdot\vect x).
\]

两个矢量并排，中间什么符号都不放。这不是偷懒，而是有道理的：$\vect a\cdot\vect b$ 是数、$\vect a\times\vect b$ 是矢量、$\vect a\vect b$ 是张量，\textbf{中间那个符号正好标记了结果的类型}。空着，就表示“既不缩并也不叉”。

\textbf{什么时候会看混。}当 $\vect a$、$\vect b$ 是带模长的符号（如 $a$ 表示长度）时，$ab$ 会被读成两个数相乘。处理办法只有一条，而且必须坚持：\textbf{看上下文里 $ab$ 后面有没有再作用在矢量上。}有，就是并矢；没有，就是数。

本讲义的约定：
\begin{itemize}
\item 正文一律写 $\vect a\otimes\vect b$，与课件 p9 一致，不会看混；
\item 但你读 Gibbs 的书、读 Landau、读连续介质教材时，看到的 $\vect a\vect b$ 就是同一个东西，直接当 $\otimes$ 处理；
\item §6 的字典里三种写法并排列出，随时可以对照。
\end{itemize}

\subhdr{5.2\quad 左右作用：并矢最容易错的地方}

并矢有两套作用规则，方向相反：

\[
\boxed{\;(\vect a\otimes\vect b)\cdot\vect x=\vect a\,(\vect b\cdot\vect x)\;}
\qquad\text{右作用：先投到 }\vect b\text{，再沿 }\vect a\text{ 放大}
\]
\[
\boxed{\;\vect x\cdot(\vect a\otimes\vect b)=(\vect x\cdot\vect a)\,\vect b\;}
\qquad\text{左作用：先投到 }\vect a\text{，再沿 }\vect b\text{ 放大}
\]

\textbf{两者一般不同。}取 $\vect a=\vect e_1$、$\vect b=\vect e_2$、$\vect x=\vect e_1+2\vect e_2+3\vect e_3$：

\[
(\vect e_1\otimes\vect e_2)\cdot\vect x=\vect e_1(\vect e_2\cdot\vect x)=2\vect e_1,
\qquad
\vect x\cdot(\vect e_1\otimes\vect e_2)=(\vect x\cdot\vect e_1)\vect e_2=\vect e_2 .
\]

一个指向 $\vect e_1$，一个指向 $\vect e_2$，差别一眼可见。

\textbf{什么时候两者相同：当并矢对称时。}力学里的 $\ten J$、$\ten\sigma$ 都对称，所以“碰巧”怎么写都对；但课件 p10 的旋转算子 $S=\sum_k\vect n'_k\otimes\vect n_k$ \textbf{不对称}，左右一混立刻出错。

\begin{center}
\fcolorbox{violet}{bandviolet}{\parbox{0.94\linewidth}{\centering\vspace{3pt}
\small 一句话：$\vect a\otimes\vect b$ 不对称，$\vect a\otimes\vect b+\vect b\otimes\vect a$ 才对称。并矢带着顺序信息，所以它不是“两个矢量的乘积”。\vspace{3pt}}}
\end{center}

\subhdr{5.3\quad 点乘、双点积、迹}

两个并矢之间可以“乘”，乘掉几个指标就叫几点积：

\[
(\vect a\otimes\vect b)\cdot(\vect c\otimes\vect d)=(\vect b\cdot\vect c)\,\vect a\otimes\vect d
\qquad\text{（缩并掉一个指标，剩二阶）},
\]
\[
(\vect a\otimes\vect b):(\vect c\otimes\vect d)=(\vect a\cdot\vect c)(\vect b\cdot\vect d)
\qquad\text{（缩并掉两个指标，剩标量）}.
\]

第二条叫\textbf{双点积}，写成指标就是 $\ten A:\ten B=\sum_{i,j}A_{ij}B_{ij}$，与矩阵写法 $\operatorname{tr}(\ten A^{\mathsf T}\ten B)$ 等价。

\textbf{迹}：$\operatorname{tr}\ten D=\sum_i D_{ii}$。用基并矢写，$\operatorname{tr}(\vect a\otimes\vect b)=\vect a\cdot\vect b$；于是

\[
\unitd:\ten D=\operatorname{tr}\ten D,\qquad
\ten A:\vect u\otimes\vect v=\vect u\cdot\ten A\cdot\vect v .
\]

\textbf{数字例子。}取 $\ten D=\vect e_1\otimes\vect e_1+2\vect e_1\otimes\vect e_2+3\vect e_2\otimes\vect e_2$，矩阵是 $\begin{pmatrix}1&2\\0&3\end{pmatrix}$（二维示意，三维同理）。则

\[
\operatorname{tr}\ten D=1+3=4,\qquad
\ten D:\ten D=1^2+2^2+3^2=14,\qquad
\unitd:\ten D=4 .
\]

最后一条是“$\unitd$ 与谁双点积就取谁的迹”，可以当公式用。

\subhdr{5.4\quad 单位并矢与投影并矢}

\textbf{单位并矢}就是克罗内克符号那个张量：

\[
\unitd=\sum_{k=1}^{3}\vect e_k\otimes\vect e_k
\quad\Longrightarrow\quad
\unitd\cdot\vect x=\sum_k\vect e_k(\vect e_k\cdot\vect x)=\vect x
\qquad\text{（对任意正交基成立，§4.4 已证）}.
\]

\textbf{投影并矢。}设 $\vect n$ 是单位矢量，令 $\ten P=\vect n\otimes\vect n$，则

\[
\ten P\cdot\vect x=\vect n(\vect n\cdot\vect x)
\quad\text{（$\vect x$ 在 }\vect n\text{ 上的投影）},
\qquad
\ten P^2=\ten P\ \text{（再投一次不变）}.
\]

而 $\unitd-\vect n\otimes\vect n$ 是“投影到垂直于 $\vect n$ 的平面”的算子。两者的迹分别是 $1$ 与 $2$，正对应“投影掉几个自由度”。
课件 p11 里那个绕 $\vect n$ 转 $\varphi$ 的旋转算子就是这两个投影并矢拼起来的：

\[
S(\varphi;\vect n)=\vect n\otimes\vect n+(\unitd-\vect n\otimes\vect n)\cos\varphi+\vect n^{\times}\sin\varphi .
\]

读法：\textbf{沿轴的分量不动（第一项），垂直轴的分量被转了 $\cos$ 又在 $\vect n\times$ 方向上补了 $\sin$（后两项）。}

\subhdr{5.5\quad 转置、对称分解、对偶张量}

\[
(\vect a\otimes\vect b)^{\mathsf T}=\vect b\otimes\vect a,
\qquad
(\ten A\cdot\ten B)^{\mathsf T}=\ten B^{\mathsf T}\cdot\ten A^{\mathsf T},
\qquad
(\ten D^{\mathsf T})_{ij}=D_{ji}.
\]

转置的判别式很好用：\textbf{对任意 $\vect x,\vect y$ 有 $(\ten D\cdot\vect x)\cdot\vect y=\vect x\cdot(\ten D^{\mathsf T}\cdot\vect y)$}。§3.5 证特征向量正交时就用的这一条。

\textbf{对称分解。}任何并矢都能唯一拆成对称与反对称两半：

\[
\ten D=\underbrace{\tfrac12(\ten D+\ten D^{\mathsf T})}_{\text{对称}}\;+\;
\underbrace{\tfrac12(\ten D-\ten D^{\mathsf T})}_{\text{反对称}} .
\]

力学里这两半各有明确身份：\textbf{对称半}给惯量、应力、应变；\textbf{反对称半}等价于一个矢量，正是叉乘。

\textbf{反对称张量与对偶张量。}若 $\ten A$ 反对称，则存在唯一的矢量 $\vect a$ 使

\[
\ten A\cdot\vect x=\vect a\times\vect x
\qquad
\left(A_{ij}=-\varepsilon_{ijk}a_k,\quad a_k=-\tfrac12\varepsilon_{kij}A_{ij}\right).
\]

这个 $\ten A$ 记作 $\vect a^{\times}$，就是课件 p8 的对偶张量，矩阵是

\[
\vect a^{\times}=
\begin{pmatrix}
0&-a_3&a_2\\
a_3&0&-a_1\\
-a_2&a_1&0
\end{pmatrix}.
\]

\textbf{这句话能省很多困惑：叉乘不是“另一种乘法”，而是一个反对称并矢在作用。}
“为什么叉乘不满足结合律”就通了：不满足结合律的是叉乘这个运算，而反对称张量的\textbf{合成}（矩阵乘法）是满足结合律的。
课件 p12 从 $S S^{\mathsf T}=\unitd$ 出发得到 $\dot S S^{\mathsf T}$ 反对称、进而得到角速度 $\vect\Omega^{\times}=\dot S S^{\mathsf T}$，走的正是这条路。

\aside{课件 p8 写的是 $a^{\times}_{ik}=\varepsilon_{ijk}a_j$，本讲义写 $A_{ij}=-\varepsilon_{ijk}a_k$。两者是\textbf{同一个矩阵}，只差改写哑指标：把后者的 $j,k$ 互换并用 $\varepsilon_{ijk}=-\varepsilon_{ikj}$ 补负号，即得前者。用分量核一遍最省事：第三种写法给出 $A_{23}=-\varepsilon_{231}a_1=-a_1$，与课件矩阵的 $(2,3)$ 元一致。}

\subhdr{5.6\quad 缩并层级表：一张能查的表}

这张表只是把 §5.3--§5.5 的规则按“缩并掉几个指标”读了一遍。

\begin{center}
\tabtop
\begin{tabular}{@{}lll@{}}
\thead{写法} & \thead{结果} & \thead{分量} \\
\multicolumn{3}{@{}l@{}}{\color{petal}\rule{\linewidth}{0.5pt}}\\
$\ten A\cdot\ten B$ & 二阶张量 & $A_{ik}B_{kj}$ \\
$\ten A\cdot\vect x$ & 矢量 & $A_{ij}x_j$ \\
$\vect x\cdot\ten A$ & 矢量 & $x_iA_{ij}$ \\
$\ten A:\ten B$ & 标量 & $A_{ij}B_{ij}$ \\
$\operatorname{tr}\ten A$ & 标量 & $A_{ii}$ \\
$(\vect a\otimes\vect b)\cdot(\vect c\otimes\vect d)$ & 二阶张量 & $(\vect b\cdot\vect c)\,\vect a\otimes\vect d$ \\
$\vect a\cdot(\vect b\otimes\vect c)$ & 矢量 & $(\vect a\cdot\vect b)\vect c$ \\
\end{tabular}
\tabbot
\end{center}

逐行读法：前五行的规律是“每点一下吃掉一个指标”，所以二阶$\cdot$二阶留二阶、二阶$\cdot$矢量留一阶、双点积吃掉两个留零阶。
最后两行要小心，它们是整行最关键的地方：

\begin{itemize}
\item 第六行：两个并矢点乘，缩并的是\textbf{右并矢的左因子}与\textbf{左并矢的右因子}，即 $\vect b$ 与 $\vect c$。这一条是课件 p11 证明 $S^{\mathsf T}S=\unitd$ 的关键一步：
$S^{\mathsf T}S=\sum_{k,l}(\vect n_k\otimes\vect n'_k)\cdot(\vect n'_l\otimes\vect n_l)=\sum_{k,l}(\vect n'_k\cdot\vect n'_l)\,\vect n_k\otimes\vect n_l=\sum_k\vect n_k\otimes\vect n_k=\unitd$。
\item 第七行：矢量点并矢，$\vect a\cdot(\vect b\otimes\vect c)=(\vect a\cdot\vect b)\vect c$，结果是矢量，方向沿 $\vect c$。它和第七行差在一个“谁先被缩并”，别凭形状猜。
\end{itemize}

\subhdr{5.7\quad 谱分解：把并矢搬到主轴上}

§3.5 的谱定理用并矢写出来，就是一句话：

\[
\ten D=\lambda_1\,\vect n_1\otimes\vect n_1+\lambda_2\,\vect n_2\otimes\vect n_2+\lambda_3\,\vect n_3\otimes\vect n_3 .
\]

这三项每一项都是一个\textbf{投影并矢}乘上系数，所以“沿主轴分解”的物理含义是：\textbf{把任意矢量先投影到三个主轴上，各自乘 $\lambda_k$，再拼回来。}

\textbf{椭球图像（接 §3.3）。}对对称正定的 $\ten D$，$\vect x\cdot\ten D\cdot\vect x=1$ 是椭球，三个半轴 $1/\sqrt{\lambda_k}$，方向就是主轴。惯量椭球、应力椭球都是这个图像。
若某个 $\lambda=0$，椭球退化成柱面或平面——这意味着那台机器在某个方向上“什么都不给”，$\det\ten D=0$（§2.2）。

\textbf{三个不变量（接 §3.6）}在并矢语言里同样成立：$I_1=\operatorname{tr}\ten D$、$I_3=\det\ten D$、$I_2=\tfrac12[(\operatorname{tr}\ten D)^2-\operatorname{tr}\ten D^2]$，做题时用来验算。

\hdr{§6}{三记号字典：并矢、指标、矩阵}

下表把同一个对象在三种写法下并排列出。\textbf{它不是要你背，而是你读到任何一本书时可以横向对照。}

\begin{center}
\begin{minipage}{\linewidth}
\small
\renewcommand{\arraystretch}{1.12}
\tabtop
\begin{tabular}{@{}p{0.30\linewidth}p{0.20\linewidth}p{0.30\linewidth}@{}}
\thead{Gibbs 并矢} & \thead{指标} & \thead{矩阵} \\
\multicolumn{3}{@{}l@{}}{\color{petal}\rule{\linewidth}{0.5pt}}\\
$\vect a\vect b$ 或 $\vect a\otimes\vect b$ & $a_ib_j$ & $\vect a\,\vect b^{\mathsf T}$（列乘行） \\
$\ten D\cdot\vect x$ & $D_{ij}x_j$ & $\ten D\,\vect x$ \\
$\vect x\cdot\ten D$ & $x_iD_{ij}$ & $\vect x^{\mathsf T}\ten D$ \\
$\ten A\cdot\ten B$ & $A_{ik}B_{kj}$ & $\ten A\ten B$ \\
$\ten A:\ten B$ & $A_{ij}B_{ij}$ & $\operatorname{tr}(\ten A^{\mathsf T}\ten B)$ \\
$\ten D^{\mathsf T}$ & $D_{ji}$ & $\ten D^{\mathsf T}$ \\
$\unitd$ & $\delta_{ij}$ & 单位矩阵 \\
$\operatorname{tr}\ten D$ & $D_{ii}$ & 对角元之和 \\
$\vect a^{\times}$ & $-\varepsilon_{ijk}a_k$ & 反对称矩阵（§5.5） \\
$\ten D=\sum_k\lambda_k\vect n_k\vect n_k$ & $\sum_k\lambda_k n_{ki}n_{kj}$ & $\ten R\operatorname{diag}(\lambda)\ten R^{\mathsf T}$ \\
\end{tabular}
\tabbot
\end{minipage}
\end{center}

\textbf{表的最后一行值得单独看。}它就是“对角化”的全部内容：找一组正交基，使矩阵变成对角阵。
矩阵写法要写三个东西（$\ten R$、对角线、$\ten R^{\mathsf T}$），并矢写法只写一项求和，而物理含义在并矢写法里是\textbf{看得见}的：三个投影并矢各自乘上主值。这就是为什么力学书偏爱并矢。

\textbf{但并矢不是矩阵。}矩阵是“坐标下的账本”，并矢是不依赖坐标系的对象。换基时账本变、对象不变（§1.3 的主线）。你可以在任意基下把并矢写成九个分量，但并矢本身不是那九个数——§4.4 已经说明，只有按变换律变的九个数才配叫张量。

\hdr{§7}{落点一：转动惯量张量}

到这里工具已经齐了。从这一节起，每个结论都只用前面出现过的东西。

\subhdr{7.1\quad 从 $\vect L=\int\vect r\times\vect v\dd m$ 推出来}

对刚体，绕定点（或取质心为基点）转动时，任一点的速度是 $\vect v=\vect\omega\times\vect r$。于是角动量

\[
\vect L=\int\vect r\times\vect v\dd m=\int\vect r\times(\vect\omega\times\vect r)\dd m .
\]

现在用“三矢量叉乘公式” $\vect a\times(\vect b\times\vect c)=\vect b(\vect a\cdot\vect c)-\vect c(\vect a\cdot\vect b)$ 展开内层：

\[
\vect r\times(\vect\omega\times\vect r)=\vect\omega(\vect r\cdot\vect r)-\vect r(\vect r\cdot\vect\omega)
=\bigl(r^2\unitd-\vect r\otimes\vect r\bigr)\cdot\vect\omega .
\]

第二步是把两项都写成“某个张量作用在 $\vect\omega$ 上”：$r^2\unitd\cdot\vect\omega=r^2\vect\omega$，而 $(\vect r\otimes\vect r)\cdot\vect\omega=\vect r(\vect r\cdot\vect\omega)$（§5.2 的作用规则反过来用）。
括号里对 $\vect\omega$ 是线性的，所以可以提到积分外面：

\[
\boxed{\;\vect L=\ten J\cdot\vect\omega,\qquad
\ten J=\int\bigl(r^2\unitd-\vect r\otimes\vect r\bigr)\dd m\;}
\]

\textbf{请注意它是并矢，而且是“各向同性减掉一个方向”的结构}：$r^2\unitd$ 让所有方向一视同仁，$\vect r\otimes\vect r$ 把沿 $\vect r$ 的那一份减掉。
正因为有第二项，$\vect L$ 才一般\textbf{不平行}于 $\vect\omega$（§7.5 给具体数值）。

\subhdr{7.2\quad 分量、几何意义、为什么对称}

把 $\ten J$ 逐分量写出来（用 $\delta_{ij}$ 表示单位并矢的分量）：

\[
J_{ij}=\int\bigl(r^2\delta_{ij}-x_ix_j\bigr)\dd m
\qquad\Longrightarrow\qquad
\ten J=
\begin{pmatrix}
\int(y^2+z^2)\dd m & -\int xy\dd m & -\int xz\dd m\\[3pt]
-\int xy\dd m & \int(x^2+z^2)\dd m & -\int yz\dd m\\[3pt]
-\int xz\dd m & -\int yz\dd m & \int(x^2+y^2)\dd m
\end{pmatrix}.
\]

读法：对角元是\textbf{绕对应坐标轴的转动惯量}；非对角元是\textbf{惯性积的负值}。
为什么非对角元带负号？它就是 $-\vect r\otimes\vect r$ 那一项贡献的 $-x_ix_j$，没有别的理由。

\textbf{对称性是白送的}：$r^2\delta_{ij}-x_ix_j$ 对 $i,j$ 对称，所以 $J_{ij}=J_{ji}$。
\textbf{正定性}也是白送的：对任意 $\vect\omega$，

\[
\vect\omega\cdot\ten J\cdot\vect\omega=\int\bigl[r^2\omega^2-(\vect r\cdot\vect\omega)^2\bigr]\dd m\ge0,
\]

因为 $|\vect r\times\vect\omega|^2=r^2\omega^2-(\vect r\cdot\vect\omega)^2\ge0$。等号只在“所有质量都在过原点且平行 $\vect\omega$ 的直线上”时成立，刚体不会这样。
于是 §3.5 的谱定理全套可用：三个实特征值、三根正交主轴。

\textbf{动能}是同一个张量给出的二次型：

\[
T=\tfrac12\int v^2\dd m=\tfrac12\int(\vect\omega\times\vect r)^2\dd m=\tfrac12\vect\omega\cdot\ten J\cdot\vect\omega .
\]

这里出现了 §3.1 那句“课件 p9 的二次型”：$T$ 就是 $\vect\omega$ 的二次型，系数表就是 $\ten J$。

\subhdr{7.3\quad 主轴 = 惯量主轴，主值 = 主转动惯量}

在主轴系里 $\ten J=\operatorname{diag}(A,B,C)$（$A,B,C$ 就是三个主转动惯量），于是

\[
\vect L=(A\omega_1,\;B\omega_2,\;C\omega_3),
\qquad
T=\tfrac12\bigl(A\omega_1^2+B\omega_2^2+C\omega_3^2\bigr).
\]

三个主值满足三角形不等式 $A+B\ge C$（循环成立）。验证只需一行：

\[
A+B-C=\int\bigl[(y^2+z^2)+(x^2+z^2)-(x^2+y^2)\bigr]\dd m=2\int z^2\dd m\ge0 .
\]

这条不等式是很便宜的验算工具：算出来的三个主值若不满足，一定哪一步错了。

\textbf{对称轴的方向就是主轴。}如果刚体有一个几何对称轴，那么绕该轴的任意转动都保持刚体不变，因此这根轴必是主轴。这条判断能省掉大量特征值计算：

\begin{itemize}
\item 均质长方体：三根棱方向都是主轴（每根轴转 $180^\circ$ 都不变）；
\item 正圆锥、圆柱：对称轴是主轴，另两根在垂直于它的平面内任意取（那两个主值相等）；
\item 均质正四面体：三根主轴过对棱中点连线，三个主值相等。
\end{itemize}

\subhdr{7.4\quad 两个必背结果与平行轴定理}

\textbf{长方体}（边长 $a,b,c$，质量 $m$，质心为原点，棱沿坐标轴）：

\[
\ten J=\frac{m}{12}\operatorname{diag}\bigl(b^2+c^2,\;a^2+c^2,\;a^2+b^2\bigr).
\]

\textbf{矩形薄板}（$a\times b$，厚度忽略，板面为 $xy$ 面）：

\[
\ten J=\frac{m}{12}\operatorname{diag}\bigl(b^2,\;a^2,\;a^2+b^2\bigr)
\qquad\Longrightarrow\qquad J_z=J_x+J_y .
\]

最后这个等式就是\textbf{垂直轴定理}。它在并矢语言里不是新定理，只是一行读数：$z$ 方向没有厚度，所以 $z$ 的惯量恰好是另两个之和。

\textbf{平行轴定理的并矢版}（比标量版好用得多）：把原点从质心 $C$ 挪到 $O$，位移 $\vect R$，则

\[
\ten J_O=\ten J_C+m\bigl(R^2\unitd-\vect R\otimes\vect R\bigr).
\]

推导只用一条 $\int\vect\rho\dd m=\vect0$（质心定义）。\textbf{注意第二项是并矢，不能写成 $mR^2$ 一个标量}——绕不同方向的轴平移，增量不同。标量版 $J=J_C+mR^2$ 只是它在“轴平行且过 $O$”这一个特例上的读数。

\subhdr{7.5\quad 具体算一遍：$\vect L$ 到底偏多少}

取均质长方体 $a=2$、$b=3$、$c=4$ m，质量 $m=5$ kg，质心在原点。由 §7.4，

\[
\ten J=\frac{5}{12}\operatorname{diag}(9+16,\;4+16,\;4+9)=\frac{5}{12}\operatorname{diag}(25,20,13)
\approx\operatorname{diag}(10.42,\;8.33,\;5.42)\ \text{kg·m}^2 .
\]

现在让它绕体对角线方向转，$\vect\omega=\omega_0(1,1,1)/\sqrt3$，取 $\omega_0=1$ rad/s。则

\[
\vect L=\ten J\cdot\vect\omega=\frac{5}{12\sqrt3}(25,\;20,\;13)\approx(6.02,\;4.81,\;3.13)\ \text{kg·m}^2/\text{s}.
\]

$\vect L$ 与 $\vect\omega\propto(1,1,1)$ 明显不平行。夹角由

\[
\cos\theta=\frac{\vect L\cdot\vect\omega}{|\vect L|\,|\vect\omega|}
=\frac{58}{\sqrt3\,\sqrt{1194}}\approx0.9691
\quad\Longrightarrow\quad
\theta\approx14.3^\circ .
\]

\textbf{这就是“刚体转起来会晃”的数学原因}：角动量不沿转轴，要维持这种转动必须持续施加力矩（§8 会把这个力矩算出来）。

\aside{上式的分子分母：$\vect L\cdot\vect\omega=\frac{5}{12\cdot3}(25+20+13)=\frac{290}{36}$，$|\vect L|=\frac{5}{12\sqrt3}\sqrt{1194}$，相除时 $\frac{5}{12}$ 与 $\frac1{\sqrt3}$ 全部约掉，只剩 $58/(\sqrt3\sqrt{1194})$。这一条在本讲义交付前已用数值脚本复核。}

\hdr{§8}{落点二：欧拉方程与对称陀螺}

\subhdr{8.1\quad 输运定理：为什么不能直接对 $\vect L$ 求导}

初学者最常写的一步错是

\[
\frac{\dd\vect L}{\dd t}=\ten J\cdot\frac{\dd\vect\omega}{\dd t}
\qquad\text{（错，除非 }\ten J\text{ 在该系里是常量）}.
\]

$\ten J$ 是“跟着刚体转”的张量。在固定系里看，它的九个分量随时间变；在体坐标系里看，它才是不变的。两个系里的时间导数由一个公式联系：

\[
\left(\frac{\dd\ten D}{\dd t}\right)_{\text{固定}}
=
\left(\frac{\dd\ten D}{\dd t}\right)_{\text{体}}
+\vect\Omega\times\ten D-\ten D\times\vect\Omega ,
\]

其中 $\vect\Omega$ 是体坐标系相对固定系的角速度。\textbf{先用一个检验确认它没写错}：取 $\ten D=\unitd$，左边为 $\vect0$（单位张量永远是单位张量），右边

\[
\vect\Omega\times\unitd-\unitd\times\vect\Omega=\vect\Omega\times-\vect\Omega\times=\vect0 .\ \checkmark
\]

对矢量（一阶张量）同一条公式退化成熟悉的形式

\[
\left(\frac{\dd\vect A}{\dd t}\right)_{\text{固定}}
=\left(\frac{\dd\vect A}{\dd t}\right)_{\text{体}}+\vect\Omega\times\vect A ,
\]

这一条在讲相对运动时已经用过；并矢版只是把它作用到三个基并矢上再线性拼起来。

\aside{为什么并矢版多出一项：一阶只有一个"$\times\vect\Omega$"，二阶有两个指标，每个指标都要被转一次，所以是 $\vect\Omega\times\ten D-\ten D\times\vect\Omega$ 两项。取 $\ten D=\vect a\otimes\vect b$ 逐项核对即可：$(\vect\Omega\times\vect a)\otimes\vect b+\vect a\otimes(\vect\Omega\times\vect b)$ 正是这两项。}

\subhdr{8.2\quad 体坐标系里的角动量定理与欧拉方程}

固定系里 $\dd\vect L/\dd t=\vect M$。换成体坐标系（体坐标系以 $\vect\omega$ 转动），用一阶的输运公式：

\[
\left(\frac{\dd\vect L}{\dd t}\right)_{\text{固定}}
=\left(\frac{\dd\vect L}{\dd t}\right)_{\text{体}}+\vect\omega\times\vect L=\vect M .
\]

体坐标系里 $\ten J=\operatorname{diag}(A,B,C)$ 是常量，$\vect L=(A\omega_1,B\omega_2,C\omega_3)$，于是 $(\dd\vect L/\dd t)_{\text{体}}=(A\dot\omega_1,B\dot\omega_2,C\dot\omega_3)$。再算 $\vect\omega\times\vect L$ 的三个分量：

\[
(\vect\omega\times\vect L)_1=\omega_2(C\omega_3)-\omega_3(B\omega_2)=(C-B)\omega_2\omega_3,
\]

另两个循环。合起来就是\textbf{欧拉方程}：

\[
A\dot\omega_1+(C-B)\omega_2\omega_3=M_1,\quad
B\dot\omega_2+(A-C)\omega_3\omega_1=M_2,\quad
C\dot\omega_3+(B-A)\omega_1\omega_2=M_3 .
\]

无外力矩（$M_k=0$）时，这三条就是刚体自由转动的全部动力学。

\subhdr{8.3\quad 对称陀螺：$A=B\ne C$}

取 $A=B$、$\vect M=\vect0$。第三条方程立刻给 $C\dot\omega_3=0$，即 $\omega_3=\text{const}$。前两条变成

\[
\dot\omega_1=-\Omega\,\omega_2,\qquad \dot\omega_2=\Omega\,\omega_1,
\qquad
\Omega\equiv\frac{C-A}{A}\,\omega_3 .
\]

这是标准的“振子方程”，解为

\[
\omega_1=\omega_\perp\cos(\Omega t+\varphi_0),\qquad
\omega_2=\omega_\perp\sin(\Omega t+\varphi_0),
\]

即\textbf{角速度矢量的横向分量绕体轴 $3$ 匀速转动}，转速 $\Omega$、大小 $\omega_\perp$ 不变。
（严格说，这里是 $\vect\omega$ 在体坐标系里绕 $\vect e_3$ 转，不是陀螺本身在进动。）

\textbf{几何图像。}$\vect L=(A\omega_1,A\omega_2,C\omega_3)$。因为 $\vect L$ 的前两个分量只比 $\vect\omega$ 的多个常数因子 $A$，所以 $\vect\omega,\vect L,\vect e_3$ \textbf{三个矢量始终共面}，这个平面绕 $\vect e_3$ 以 $\Omega$ 转。这是对称陀螺最有用的图像。

\begin{center}
\fcolorbox{violet}{bandviolet}{\parbox{0.94\linewidth}{\centering\vspace{3pt}
\small $C>A$（扁盘状，如地球）时 $\Omega>0$，$\vect\omega$ 与 $\vect L$ 绕 $\vect e_3$ 同向转；$C<A$（细长杆状）时反向。\vspace{3pt}}}
\end{center}

\subhdr{8.4\quad 回到 §7.5：那个 $14.3^\circ$ 的力矩有多大}

把 $\vect L$ 与 $\vect\omega$ 不平行这件事代进 §8.2 的 $\vect M=\vect\omega\times\vect L$，就得到“维持这种转动需要多大扭矩”。

对 §7.5 的长方体，若强行让 $\vect\omega$ 沿体对角线（即强行让它不晃），所需力矩为 $\vect M=\vect\omega\times\vect L$，其大小

\[
|\vect M|=\omega_0\,|\vect L|\sin\theta
\]

不为零，方向垂直于 $\vect\omega$ 与 $\vect L$ 张成的平面。\textbf{自由转动时没有这个力矩，所以 $\vect\omega$ 会自己动起来}——这正是 §8.3 那条进动的来源。
一句话：\textbf{角动量不平行于角速度，就必须有外力矩；没有外力矩，角速度就得改。}

\hdr{§9}{落点三：应力张量}

\subhdr{9.1\quad 面力是法向的函数：商法则的用武之地}

在连续体里过一点取一个截面，法向 $\vect n$，截面上单位面积的受力记 $\vect t(\vect n)$，叫\textbf{面力矢量}。
柯西应力定理说：$\vect t$ 是 $\vect n$ 的\textbf{线性}函数，即存在一个二阶张量 $\ten\sigma$ 使

\[
\vect t(\vect n)=\ten\sigma\cdot\vect n,
\qquad
\sigma_{ij}=\text{法向为 }\vect e_j\text{ 的面上、沿 }\vect e_i\text{ 方向的应力分量}.
\]

\textbf{为什么它一定是张量？}因为 $\vect t$ 与 $\vect n$ 都是矢量，而“对任意 $\vect n$ 成立 $\vect t=\ten\sigma\vect n$”正是 §4.5 商法则的前提，所以 $\ten\sigma$ 自动是二阶张量。
这一步不是假设，是被逼出来的——这也是本讲义 §4.5 存在的理由。

\subhdr{9.2\quad 对称性来自角动量守恒}

对微元取力矩平衡（不设体力矩），三个切应力分量必须两两相等：

\[
\sigma_{ij}=\sigma_{ji}.
\]

所以应力张量也\textbf{对称}，§3.5 的谱定理同样适用：三个主应力 $\sigma_1,\sigma_2,\sigma_3$、三个互相垂直的主平面（这些截面上剪应力为零）。
\textbf{三个主应力之和 $\sigma_1+\sigma_2+\sigma_3=\operatorname{tr}\ten\sigma$ 是与坐标无关的}——这就是“静水压强”那种不随截面变化的成分，$p=-\tfrac13\operatorname{tr}\ten\sigma$。

\subhdr{9.3\quad 单轴拉伸：一个能口算的例子}

取 $\ten\sigma=\sigma_0\vect e_1\otimes\vect e_1$（只有 $\sigma_{11}=\sigma_0$）。主轴就是 $\vect e_1,\vect e_2,\vect e_3$，主应力 $\sigma_0,0,0$。

在 $xy$ 面内取 $45^\circ$ 截面，$\vect n=(\vect e_1+\vect e_2)/\sqrt2$：

\[
\vect t=\ten\sigma\cdot\vect n=\sigma_0\bigl[\vect e_1(\vect e_1\cdot\vect n)\bigr]=\frac{\sigma_0}{\sqrt2}\vect e_1 .
\]

法向分量 $\vect t\cdot\vect n=\sigma_0/2$，切向分量的大小 $\sqrt{|\vect t|^2-(\sigma_0/2)^2}=\sigma_0/2$。
\textbf{最大剪应力出现在与主轴成 $45^\circ$ 的截面上，等于 $\sigma_0/2$}——这就是“拉伸试件沿 $45^\circ$ 断裂”的并矢解释。它不需要背，只要把 $\ten\sigma\cdot\vect n$ 算一遍。

\subhdr{9.4\quad 平衡方程：高斯定理的并矢版}

把面力沿闭合曲面加起来，用高斯定理换成体积分：

\[
\oint_S\ten\sigma\cdot\vect n\dd S=\int_V\nabla\cdot\ten\sigma\dd V ,
\qquad
(\nabla\cdot\ten\sigma)_k=\sum_i\partial_i\sigma_{ik} .
\]

与体力 $\vect f$ 平衡得

\[
\nabla\cdot\ten\sigma+\vect f=\vect0\quad\text{（静力学）},
\qquad
\nabla\cdot\ten\sigma+\vect f=\rho\vect a\quad\text{（动力学）}.
\]

写得开就是 $\partial_j\sigma_{ji}+f_i=0$——三个方程，正好对应三个方向的力平衡。

\warnbox{\textbf{两个容易混的记号。}$\nabla\cdot\ten\sigma$ 是\textbf{矢量}（对第二个指标缩并）；$\nabla\ten\sigma$ 是\textbf{三阶}张量。两者差一个点，差一个阶。
另一个坑：缩并的是第二个指标 $\partial_j\sigma_{ji}$，不是第一个——这跟 $\ten\sigma$ 的指标定义（第一个指标是“沿哪个方向受力”）直接相关，写反了在非对称情形会出错。}

\hdr{§10}{四道例题}

本节四道题都是\textbf{自编例题}，不来自课件习题——课件 p16--p18 的 E8.01、E8.02 与 HW8.01--HW8.03 全是运动学题（点的速度、瞬心），已在《Topic 8 刚体运动学 P1 精讲》里逐题处理；本讲义练的是张量与并矢的运算，故另配题目。每题末尾标出它用到本讲义哪一节。

\qbreak
\qstem{例 1（自编）\quad 设 $\ten D=4\vect e_1\otimes\vect e_1+\vect e_1\otimes\vect e_2+\vect e_2\otimes\vect e_1+4\vect e_2\otimes\vect e_2+2\vect e_3\otimes\vect e_3$。求 $\ten D\cdot\vect x$（$\vect x=\vect e_1+2\vect e_2+3\vect e_3$）、$\operatorname{tr}\ten D$、$\det\ten D$、$\ten D:\ten D$、特征值与主轴。}

\textbf{解。}先把九个系数按行列摆好（第 $i$ 行第 $j$ 列是 $\vect e_i\otimes\vect e_j$ 的系数）：

\[
\ten D=\begin{pmatrix}4&1&0\\1&4&0\\0&0&2\end{pmatrix}.
\]

作用：$(\ten D\vect x)_i=\sum_j D_{ij}x_j$，第一行 $4\cdot1+1\cdot2+0=6$，第二行 $1\cdot1+4\cdot2=9$，第三行 $0+0+2\cdot3=6$，所以

\[
\ten D\cdot\vect x=(6,\;9,\;6).
\]

因为 $\ten D$ 对称，左作用 $\vect x\cdot\ten D$ 给同一个结果（§5.2）。

迹与行列式：$\operatorname{tr}\ten D=4+4+2=10$；$\det\ten D=2(16-1)=30$（按第三行展开，前两个 $2\times2$ 子式只有一个非零）。
双点积：$\ten D:\ten D=\sum_{i,j}D_{ij}^2=16+1+1+16+4=38$。

特征方程由 §2.3：$-\lambda^3+10\lambda^2-I_2\lambda+30=0$，其中 $I_2=\tfrac12(10^2-38)=31$。
即 $\lambda^3-10\lambda^2+31\lambda-30=0$，试根 $\lambda=2$：$8-40+62-30=0$，成立。分解得 $\lambda=2,3,5$。

主轴：$\lambda=5$ 时 $(\ten D-5\unitd)$ 的前两块给出 $-\vect e_1+\vect e_2$ 方向，即 $\vect n_1=(\vect e_1+\vect e_2)/\sqrt2$；
$\lambda=3$ 时得 $\vect n_2=(-\vect e_1+\vect e_2)/\sqrt2$；$\lambda=2$ 时对应 $\vect n_3=\vect e_3$。谱分解

\[
\ten D=5\,\vect n_1\otimes\vect n_1+3\,\vect n_2\otimes\vect n_2+2\,\vect e_3\otimes\vect e_3 .
\]

验算（§3.6）：$5+3+2=10=\operatorname{tr}$，$5\cdot3\cdot2=30=\det$，都对。\checkmark

\qbreak
\qstem{例 2\quad 均匀长方体，质量 $m=5$ kg，边长 $a=2,b=3,c=4$（米），质心在原点。求 $\ten J$；再把它压成薄板（$c\to0$，质量仍为 $m$），验证垂直轴定理。}

\textbf{解。}由 §7.4 直接代入：

\[
\ten J=\frac{m}{12}\operatorname{diag}(b^2+c^2,\;a^2+c^2,\;a^2+b^2)
=\frac{5}{12}\operatorname{diag}(25,\;20,\;13)
\approx\operatorname{diag}(10.42,\;8.33,\;5.42)\ \text{kg·m}^2 .
\]

验算三角形不等式：$10.42+8.33=18.75\ge5.42$，其余两个循环也成立（§7.3）。

压成薄板（$c=0$）后

\[
\ten J_{\text{板}}=\frac{m}{12}\operatorname{diag}(b^2,\;a^2,\;a^2+b^2)
=\frac{5}{12}\operatorname{diag}(9,\;4,\;13),
\]

于是 $J_z=13m/12$，而 $J_x+J_y=(9+4)m/12=13m/12$，二者相等，垂直轴定理成立。\checkmark

\qbreak
\qstem{例 3\quad 对称陀螺 $A=B=2$ kg·m$^2$，$C=3$ kg·m$^2$，自转 $\omega_3=10$ rad/s，无外力矩。求横向分量的进动角速度与周期。}

\textbf{解。}由 §8.3，

\[
\Omega=\frac{C-A}{A}\omega_3=\frac{3-2}{2}\times10=5\ \text{rad/s},
\qquad
T=\frac{2\pi}{\Omega}=\frac{2\pi}{5}\approx1.26\ \text{s}.
\]

即 $\vect\omega$ 的横向分量每 $1.26$ 秒绕体轴转一圈，$\omega_3$ 与 $|\vect\omega|$ 都不变。
$C>A$，属于“扁盘状”，$\vect\omega$ 与 $\vect L$ 同向绕 $\vect e_3$ 转（§8.3 的框）。

\qbreak
\qstem{例 4\quad 仍用例 2 的 $\ten J=\operatorname{diag}(10.42,8.33,5.42)$。求绕体对角线方向 $\vect n=(1,1,1)/\sqrt3$ 的转动惯量，并验证 $\vect\omega\cdot\vect L=2T$。}

\textbf{解。}绕单位方向 $\vect n$ 的转动惯量是二次型在 $\vect n$ 上的值（瑞利商，§3.5 的旁注）：

\[
I(\vect n)=\vect n\cdot\ten J\cdot\vect n
=\frac{1}{3}(10.42+8.33+5.42)=\frac{24.17}{3}\approx8.06\ \text{kg·m}^2 .
\]

它落在最小主值 $5.42$ 与最大主值 $10.42$ 之间，符合 $\lambda_{\min}\le I\le\lambda_{\max}$。

验证：取 $\vect\omega=\vect n$（$|\vect\omega|=1$），则 $T=\tfrac12\vect\omega\cdot\ten J\cdot\vect\omega=\tfrac12 I=4.03$；
而 $\vect L=\ten J\cdot\vect\omega=\frac{1}{\sqrt3}(10.42,8.33,5.42)$，于是

\[
\vect\omega\cdot\vect L=\frac{1}{\sqrt3}\cdot\frac{1}{\sqrt3}(10.42+8.33+5.42)=8.06=2T .\ \checkmark
\]

注意 $\vect\omega\cdot\vect L=2T$ 是恒等式（不依赖主轴），它说明“功率 $=\vect\omega\cdot\vect M$”那条与这里的二次型是同一件事。

\hdr{§11}{易错点与速查}

\warnbox{\textbf{七个高频错误}\\[2pt]
1. \textbf{把张量当成“九个数字”}：九个数字要按变换律 $T'_{ij}=R_{ki}R_{lj}T_{kl}$ 变才叫张量（§4.4）。不检验就套用谱定理，是错的源头。\\
2. \textbf{把 $\vect a\otimes\vect b$ 当成 $\vect a\cdot\vect b$}：前者是算子（九个分量），后者是标量。中间是 $\otimes$ 还是 $\cdot$，差一个层级。\\
3. \textbf{左右的顺序混用}：$\ten D\cdot\vect x\ne\vect x\cdot\ten D$，除非 $\ten D$ 对称。力学里 $\ten J$、$\ten\sigma$ 对称，所以“碰巧”不错；换到旋转算子 $S$ 上立刻出错（§5.2）。\\
4. \textbf{直接写 $\dd\vect L/\dd t=\ten J\cdot\dot{\vect\omega}$}：只在 $\ten J$ 于该系为常量时成立；体坐标系里求的是体导数，固定系里还要补 $\vect\omega\times\ten J\cdot\vect\omega$（§8.1）。\\
5. \textbf{把惯量张量的非对角元符号写反}：$J_{12}=-\int xy\dd m$，负号不能丢；丢了以后主轴会算错。\\
6. \textbf{把 $\nabla\cdot\ten\sigma$ 当成标量或三阶张量}：它是对第二个指标缩并出来的\textbf{矢量}（§9.4）。\\
7. \textbf{把特征值当成“唯一的解”}：三个特征值不变，但\textbf{特征向量的分量}随基变（§3.6）。说“主轴方向”必须连着坐标系说。}

\begin{center}
\begin{minipage}{\linewidth}
\tabtop
\begin{tabular}{@{}p{0.34\linewidth}p{0.56\linewidth}@{}}
\thead{要做的事} & \thead{用什么} \\
\multicolumn{2}{@{}l@{}}{\color{petal}\rule{\linewidth}{0.5pt}}\\
判断是不是张量 & 查变换律 $T'_{ij}=R_{ki}R_{lj}T_{kl}$；或对已知矢量的线性关系用商法则（§4.5）\\[3pt]
把“矢量 $\to$ 矢量”写成式子 & $\ten D=D_{ij}\vect e_i\otimes\vect e_j$，作用规则 $(\vect a\otimes\vect b)\cdot\vect x=\vect a(\vect b\cdot\vect x)$\\[3pt]
判断是否对称 & 看 $D_{ij}=D_{ji}$；反对称部分等价于轴矢量（叉乘）\\[3pt]
求主轴与主值 & 解 $\det(\ten D-\lambda\unitd)=0$；对称时主轴正交，$\ten D=\sum\lambda_k\vect n_k\otimes\vect n_k$\\[3pt]
验算 & $I_1=\operatorname{tr}$、$I_3=\det$；惯量张量还有 $A+B\ge C$\\[3pt]
转动惯量 & $\ten J=\int(r^2\unitd-\vect r\otimes\vect r)\dd m$；平行轴 $\ten J_O=\ten J_C+m(R^2\unitd-\vect R\otimes\vect R)$\\[3pt]
角动量与动能 & $\vect L=\ten J\cdot\vect\omega$，$T=\frac12\vect\omega\cdot\ten J\cdot\vect\omega$，$\vect\omega\cdot\vect L=2T$\\[3pt]
体坐标系动力学 & $(\dd\vect L/\dd t)_{\text{体}}+\vect\omega\times\vect L=\vect M$，即欧拉方程\\[3pt]
面力 & $\vect t(\vect n)=\ten\sigma\cdot\vect n$；$\sigma_{ij}=\sigma_{ji}$；主应力 = 特征值\\
\end{tabular}
\tabbot
\end{minipage}
\end{center}

\hdr{附录 A}{指标记法速查}

\begin{center}
\begin{minipage}{\linewidth}
\small
\renewcommand{\arraystretch}{1.05}
\tabtop
\begin{tabular}{@{}ll@{}}
\thead{对象} & \thead{指标写法} \\
\multicolumn{2}{@{}l@{}}{\color{petal}\rule{0.55\linewidth}{0.5pt}}\\
并矢 & $D_{ij}=a_ib_j$ \\
右作用 & $(\ten D\cdot\vect x)_i=D_{ij}x_j$ \\
左作用 & $(\vect x\cdot\ten D)_j=x_iD_{ij}$ \\
张量乘张量 & $(\ten A\cdot\ten B)_{ij}=A_{ik}B_{kj}$ \\
双点积 & $\ten A:\ten B=A_{ij}B_{ij}$ \\
迹 & $\operatorname{tr}\ten D=D_{ii}$ \\
二次型 & $E=a_iD_{ij}a_j$ \\
变换律 & $D'_{ij}=R_{ki}R_{lj}D_{kl}$ \\
叉乘 & $(\vect a\times\vect b)_i=\varepsilon_{ijk}a_jb_k$ \\
对偶张量 & $(a^{\times})_{ij}=-\varepsilon_{ijk}a_k$ \\
转动惯量 & $J_{ij}=\int(r^2\delta_{ij}-x_ix_j)\dd m$ \\
散度 & $(\nabla\cdot\ten D)_k=\partial_iD_{ik}$ \\
\end{tabular}
\tabbot
\end{minipage}
\end{center}

\aside{爱因斯坦求和约定：同一项里重复出现两次的指标自动求和。本讲义正文写求和号，附录与 §7 之后的推导用约定，避免“缩并了哪个指标”看错。}

\hdr{附录 B}{自测六问}

\begin{enumerate}[label=(\arabic*)]
\item $\vect a\otimes\vect b\cdot\vect c$ 与 $\vect a(\vect b\cdot\vect c)$ 是否相同？那 $\vect a\otimes\vect b\cdot\vect c$ 与 $(\vect a\cdot\vect b)\vect c$ 呢？
\item 给定九个排成方阵的数，怎样判断它是不是张量？举一个“在旋转下像矢量、但坐标系允许平移时就不像矢量”的例子。
\item 迹为什么换基不变？请用 $\operatorname{tr}(\ten A\ten B)=\operatorname{tr}(\ten B\ten A)$ 写出两步。
\item 二次型 $\vect x\cdot\ten D\cdot\vect x$ 为什么只与 $\ten D$ 的对称部分有关？举一个 $2\times2$ 的数字例子验证。
\item $J_{12}=-\int xy\dd m$ 里的负号从哪来？把它丢掉以后，谱定理还成立吗？算出的主转动惯量还对吗？
\item 单轴拉伸 $\ten\sigma=\sigma_0\vect e_1\otimes\vect e_1$，最大剪应力多大、出现在哪个截面？为什么它不需要背？
\end{enumerate}

\aside{答案：(1) 相同，都等于“$\vect b\cdot\vect c$ 乘 $\vect a$”；但 $\vect a\otimes\vect b\cdot\vect c\ne(\vect a\cdot\vect b)\vect c$——后者是“$\vect a\cdot\vect b$ 乘 $\vect c$”，方向完全不同，见 §5.6 的第七行。
(2) 按 $T'_{ij}=R_{ki}R_{lj}T_{kl}$ 检验；位置矢量 $(x_1,x_2,x_3)$ 在旋转下合格，但坐标系平移时多出常数项，不再齐次（§4.4）。
(3) $\operatorname{tr}(\ten R^{\mathsf T}\ten T\ten R)=\operatorname{tr}(\ten R\ten R^{\mathsf T}\ten T)=\operatorname{tr}\ten T$，第一步用循环性，第二步用 $\ten R\ten R^{\mathsf T}=\unitd$（§2.1）。
(4) 交叉项 $D_{12}x_1x_2+D_{21}x_2x_1=(D_{12}+D_{21})x_1x_2$，只依赖两者的和；取 $\ten D=\begin{pmatrix}0&3\\0&0\end{pmatrix}$ 得 $E=3xy$，而对称部分 $\begin{pmatrix}0&1.5\\1.5&0\end{pmatrix}$ 给的还是 $3xy$（§3.2）。
(5) 负号来自 $-\vect r\otimes\vect r$ 那一项；丢掉后矩阵仍对称，谱定理照样成立，但那是另一个张量，主轴与主值都不再是真实的惯量主轴与主转动惯量。
(6) $\sigma_0/2$，出现在与 $\vect e_1$ 成 $45^\circ$ 的截面上；它不需要背，把 $\ten\sigma\cdot\vect n$ 算一遍就出来（§9.3）。}

\flowerrule

\subhdr{三句话总结}

\begin{enumerate}[label=\arabic*.,leftmargin=1.6em]
\item \textbf{张量是对象，矩阵是账本。}换基时账本变、对象不变；迹、行列式、特征值是账本里不随基变的数。
\item \textbf{并矢只是二阶张量的另一种写法。}作用规则 $(\vect a\otimes\vect b)\cdot\vect x=\vect a(\vect b\cdot\vect x)$ 一条就够；左右不对称、双点积吃掉两个指标、谱分解是三个投影并矢相加。
\item \textbf{理论力学里三处落点用的是同一台机器。}惯量张量给出 $\vect L=\ten J\cdot\vect\omega$，欧拉方程是它在体坐标系里的求导，应力张量是“面力对法向线性”的商法则产物。
\end{enumerate}

{\small\color{inkgray} 文中数值结果均经符号/数值复核：§1.3 的换基矩阵由 $\ten R^{\mathsf T}\ten T\ten R$ 与逐项投影两条路各算一遍；§3.3 的 $45^\circ$ 化简核对到 $3u^2+v^2$；§7.5 的夹角由 $\vect L\cdot\vect\omega/(|\vect L||\vect\omega|)$ 复算；例 1 的特征值与 $\operatorname{tr}$、$\det$、双点积一致；例 4 的 $\vect\omega\cdot\vect L=2T$ 单独验算。}

\end{document}
